\label{sec:dense_method}
We consider terminal-only, phase-only, and combined dense-and-phase rewards. The combined reward uses a manually defined sequence of reach, grasp, lift, move, and place phases, with a hand-specified progress function for each phase.

At each simulator step, the tracker uses object and gripper state to compute progress for the active phase. Reaching progress increases as gripper--object distance decreases; grasping combines proximity with gripper closure and grasp state; lifting uses height gained above the initial object position; transport uses planar object--target distance; placement uses target proximity, release, and environment success. Each progress measure is normalized to the interval from zero to one. For progress $p_t$ and the phase's previous maximum $p_{\rm best}$, the rewarded increment is $\max(p_t-p_{\rm best},0)$. A completion predicate advances the tracker to the next phase and produces a phase bonus. The weighted progress increment and phase bonus are bounded and combined with terminal feedback. The native task-success verifier remains unchanged.

In the phase-only condition, the continuous progress weight is zero, leaving phase-completion bonuses and terminal feedback.

The dense construction requires a progress function, normalization, and relevant state variables for every phase. Extending it to a new multi-stage task therefore requires further task-specific reward design. Our adaptation pipeline uses this detailed feedback only to obtain the initial base policy on simple skills. Full-task training uses binary milestones and their completion predicates, so the composition stage does not require a continuous progress function for every part of the instruction.
